% All main text, references and appendices are in this file.
% Build with: latexmk -pdf -interaction=nonstopmode -halt-on-error main.tex
\documentclass[runningheads]{llncs}
\usepackage[T1]{fontenc}
\usepackage[utf8]{inputenc}
\usepackage[expansion=false]{microtype}
\usepackage{amsmath,amssymb,mathtools}
\usepackage{mathrsfs}
\usepackage{graphicx}
\usepackage{placeins}
\graphicspath{{figures/}}
\usepackage{booktabs,tabularx,array}
\usepackage[numbers,sort&compress]{natbib}
\usepackage{xurl}
\usepackage{hyperref}
\hypersetup{hidelinks,pdftitle={Paradigm Shifts: Theory Revision and Belief Sharing in Epistemic Communities},pdfauthor={Jonas Hallgren, David Hyland, Mahault Albarracin}}
\newcolumntype{Y}{>{\raggedright\arraybackslash}X}
\newcommand{\Normal}{\mathcal N}
\newcommand{\diag}{\operatorname{diag}}
\newcommand{\KL}[2]{D_{\mathrm{KL}}\!\left(#1\,\middle\|\,#2\right)}
\newcommand{\E}{\mathbb E}
\newcommand{\tr}{\operatorname{tr}}
\newcommand{\R}{\mathbb R}
\newcommand{\softmax}{\operatorname{softmax}}
\title{Paradigm Shifts: Theory Revision and Belief Sharing in Epistemic Communities}
\titlerunning{Theory Revision and Belief Sharing in Epistemic Communities}
\author{Jonas Hallgren\inst{1} \and David Hyland\inst{2} \and Mahault Albarracin\inst{3}}
\authorrunning{J. Hallgren et al.}
% Author-supplied affiliation and email details are still required.
\institute{Affiliation 1 \email{...} \and Affiliation 2 \and Affiliation 3}
\renewcommand{\topfraction}{0.9}
\renewcommand{\bottomfraction}{0.8}
\renewcommand{\textfraction}{0.08}
\renewcommand{\floatpagefraction}{0.8}
\renewcommand{\bibsection}{\section*{References}}
\begin{document}
\raggedbottom
\maketitle
\begin{abstract}
Scientific communities revise both their estimates within an explanation and their
confidence in that explanation relative to alternatives. We model these processes
with agents that maintain conditional beliefs within candidate theories and
probabilities over those theories. Agents exchange both kinds of belief, score new
observations before incorporating them, and gradually relax conditional beliefs
toward reference distributions. We formulate the update for general probability
distributions and implement it with Gaussian models of two schematic accounts of
combustion. Following a change in the observed mass signal, the community passes
through increased uncertainty before favouring the oxygen account. When access to
observations becomes selective, communication enables agents without the changed
mass measurement to revise their preferred account and improves prediction relative
to selective observation in isolation. Network structure and the placement of
fully observing agents affect revision times. In a constant-world comparison,
communicating groups select the same account while retaining different estimates
of its quantities. These results distinguish theory choice, uncertainty, and
predictive accuracy, and relate variational belief updating to the social
organization of scientific inquiry.
\keywords{Active inference \and Theory revision \and Belief sharing \and Network epistemology}
\end{abstract}

\section{Introduction}
\label{sec:introduction}
A researcher may revise an estimated quantity while retaining the theory that
gives it meaning, or shift support toward a competing explanation. How does
evidence available to some members of a scientific community affect the
commitments of others?

Network epistemology studies how communication, network structure, and the timing
of exchange affect collective inquiry \citep{balagoyal1998,zollman2010,zollman2013}.
Agents may exchange observations or combine their probability distributions
\citep{degroot1974,genest1986}. Accounts of scientific change additionally emphasize
the organization of beliefs within theories. Quine describes dependencies among
commitments \citep{quine1951}; Lakatos distinguishes a research programme's core
from the auxiliary assumptions revised around it \citep{lakatos1970}. Kuhn's
account connects sustained commitment, crisis, and changes of theory
\citep{kuhn1962}, while Laudan examines competition between research traditions
\citep{laudan1977}. These accounts motivate a distinction between learning within
an explanation and revising confidence in the explanation itself.

Bayesian epistemology treats theory choice probabilistically \citep{bovens2003}.
Coherence models represent relations among propositions \citep{thagard1989,thagard1990},
and probabilistic approaches describe the acquisition of structured explanations
\citep{kemp2008form,ullman2020}. Active-inference research has examined epistemic
communities, shared beliefs, and the costs of belief revision
\citep{albarracin2022,friston2023federated,hyland2025}, including relations between
individual and collective representations \citep{catal2024,waade2025}. Our model
combines conditional learning within several explanations with a separate
probability distribution expressing support for them.

The communication rule builds on logarithmic opinion pooling
\citep{genest1986} and distributed hypothesis testing
\citep{jadbabaie2012,molavi2018,nedic2017,lalitha2018}. Learning from uncertain
Gaussian models has also been studied in this setting \citep{hare2019gaussian}.
We use these methods to examine how within-theory accommodation and between-theory
support interact when evidence is unevenly available. The local update has a
variational free-energy interpretation, connecting it to the perceptual component
of active inference \citep{friston2017process,dacosta2020}.

Our example draws on competing release and oxygen accounts of combustion.
Historical phlogiston theories were diverse, and their replacement was a complex
process \citep{blumenthal2017,holmes2000,boantza2011}. The example uses two
standardized candidate models to study the distinction between changing quantities
and changing explanations. We examine revision after a change in observations,
communication under selective observation, and variation across contact networks.
A complementary comparison separates disagreement about a theory's quantities
from disagreement about which theory to favour.

\section{Belief representation and updating}
\label{sec:model}
\subsection{Beliefs within and between theories}
\label{sec:representation}
There are $N$ agents, indexed by $i,j\in\{1,\ldots,N\}$, and $K$ supplied candidate
theories, indexed by $k\in\{1,\ldots,K\}$. Candidate $k$ defines a state $x_k$,
which contains its measured and unmeasured quantities, and a likelihood specifying how states predict observations in a shared space. Candidate states may differ in dimension or
interpretation. Agents using the same candidate share its coordinates. At the
start of round $t$, agent $i$ represents
\begin{equation}
 Q_{i,t}(k,x_k)=q_{ik,t}\rho_{ik,t}(x_k),\qquad
 q_{ik,t}=\frac{\exp L_{ik,t}}{\sum_{k'=1}^K\exp L_{ik',t}}.
 \label{eq:state}
\end{equation}
The normalized distribution $\rho_{ik,t}$ describes uncertainty within theory
$k$, and $q_{ik,t}$ is the probability assigned to that theory. The real numbers
$L_{ik,t}$ store its log weights. Each working distribution $\rho_{ik,t}$ has a
fixed reference $\rho^0_{ik}$ toward which it can relax. The superscript $0$
denotes a reference, whereas the subscript $t=0$ denotes the initial state.
Figure~\ref{fig:representation} illustrates these objects.

\begin{figure}[!htbp]
 \centering\includegraphics[width=\linewidth]{fig01_general_representation.pdf}
 \caption{Representation and communication. (a) Private observations are evaluated
 through candidate-specific likelihoods. (b) Agent $i$ retains conditional beliefs
 $\rho_{ik}$ and model probabilities $q_{ik}$; two candidates are shown.
 (c) Contact edges permit exchange of these beliefs. Dashed lines magnify the
 selected agent. Candidate and contact graphs are schematic; time indices are
 omitted inside the agent.}
 \label{fig:representation}
\end{figure}

Rounds are indexed by $t=0,\ldots,T-1$. All agents read the state at the start of
the round, communicate, evaluate the new observation, learn, and then apply
memory. State $t+1$ includes observation $t$. Plot coordinates labelled
``Simulation step'' report $t$ and show this post-update state. All logarithms
are natural. Integrals below use the candidate's state measure; discrete states
use sums. The normalizers are assumed positive and finite
(Appendix~\ref{app:distribution}).

\subsection{Communication, prediction, and learning}
\label{sec:update}
Let $W^x_t$ and $W^m_t$ specify communication of conditional beliefs and model
weights, respectively. Their entries are nonnegative, each row sums to one, and
entry $ij$ gives sender $j$'s influence on receiver $i$. Self-influence is
included. A bar denotes the belief after communication within round $t$:
\begin{equation}
 \bar\rho_{ik,t}(x_k)\propto\prod_{j=1}^N\rho_{jk,t}(x_k)^{W^x_{ij,t}},\qquad
 \bar L_{ik,t}=(1-\alpha_m)L_{ik,t}
       +\alpha_m\sum_{j=1}^N W^m_{ij,t}L_{jk,t}.
 \label{eq:sharing}
\end{equation}
The first operation normalizes the weighted product within each candidate,
forming a logarithmic opinion pool \citep{genest1986}. The parameter
$\alpha_m\in[0,1]$ controls model-weight sharing; conditional sharing uses $W^x_t$
directly. The two pools are normalized separately. Setting $W^x_t=I_N$ disables
conditional sharing; $\alpha_m=0$ disables model-weight sharing. Here $I_n$ is
the $n$-dimensional identity matrix.

Agent $i$ evaluates datum $y_{i,t}$ using the normalized working likelihood
$\mathcal L_{ik,t}(y\mid x_k)$, which can incorporate assigned observation weights.
Its predictive log score and subsequent update are
\begin{align}
 \ell_{ik,t}&=\log\int\mathcal L_{ik,t}(y_{i,t}\mid x_k)
                                \bar\rho_{ik,t}(x_k)\,dx_k,
 \label{eq:score}\\
 \rho^+_{ik,t}(x_k)&=e^{-\ell_{ik,t}}
       \mathcal L_{ik,t}(y_{i,t}\mid x_k)\bar\rho_{ik,t}(x_k),\qquad
 L^+_{ik,t}=\bar L_{ik,t}+\ell_{ik,t}.
 \label{eq:learning}
\end{align}
The superscript $+$ denotes the learned state before memory is applied. Scoring
precedes learning from the same datum: a candidate can accommodate an observation
while receiving a low score for its earlier prediction. Likelihoods use a common
observation measure, including their normalization constants.

Conditional learning minimizes variational free energy over trial distributions
$\varrho$:
\begin{equation}
 \mathcal F_{ik,t}[\varrho]
 =\KL{\varrho}{\bar\rho_{ik,t}}
   -\E_{\varrho}\log\mathcal L_{ik,t}(y_{i,t}\mid x_k)
 =\KL{\varrho}{\rho^+_{ik,t}}-\ell_{ik,t}.
 \label{eq:free-energy}
\end{equation}
Here $D_{\mathrm{KL}}$ is the Kullback--Leibler divergence, a nonnegative measure
of difference between distributions, and $\E_{\varrho}$ denotes expectation under
$\varrho$. The first expression balances departure from the communicated prior
against observation fit; the second shows that $\rho^+_{ik,t}$ is the unrestricted
minimizer. Appendix~\ref{app:distribution} gives the normalization and joint
model-choice objective.

Memory is a geometric interpolation between the reference and learned beliefs:
\begin{equation}
 \rho_{ik,t+1}(x_k)\propto
       [\rho^0_{ik}(x_k)]^{1-\omega}[\rho^+_{ik,t}(x_k)]^{\omega},\qquad
 L_{ik,t+1}=L^+_{ik,t},\quad 0\leq\omega\leq1.
 \label{eq:memory}
\end{equation}
The retention parameter $\omega=1$ preserves the learned conditional distribution;
$\omega=0$ resets it to its reference after scoring. Model log weights retain
the accumulated predictive scores. Information from observation $t$ is shared
with neighbours in round $t+1$.

\subsection{Gaussian implementation}
\label{sec:gaussian}
The experiments use multivariate normal (Gaussian) beliefs on four-dimensional states:
$\rho_{ik,t}=\Normal(\mu_{ik,t},\Sigma_{ik,t})$.
Their precision $\Pi_{ik,t}=\Sigma_{ik,t}^{-1}$ is the inverse covariance, and
$h_{ik,t}=\Pi_{ik,t}\mu_{ik,t}$ is the information vector. These two parameters
store the distribution and make communication and learning algebraic operations.

Raw observations have the form $o_{i,t}=u_t+\epsilon_{i,t}$, where $u_t\in\R^3$
is the prescribed common signal and $\epsilon_{i,t}\sim\Normal(0,\sigma^2I_3)$
is independent across agents and rounds. The observation matrix $H$ maps each
candidate state to the same three observables. Assigned channel weights
$w_{i,t}\in[0,1]^3$ give
\begin{equation}
 \begin{gathered}
 S_{i,t}=\diag(\sqrt{w_{i,t}}),\quad y_{i,t}=S_{i,t}o_{i,t},\quad
 \widetilde H_{i,t}=S_{i,t}H,\\
 \mathcal L_{ik,t}(y\mid x_k)=\Normal(y;\widetilde H_{i,t}x_k,\sigma^2I_3).
 \end{gathered}
 \label{eq:attention}
\end{equation}
The square root is taken componentwise and $\diag$ forms a diagonal matrix.
A positive weight $w$ assigns working raw-channel variance $\sigma^2/w$; a zero
weight contributes no candidate-discriminating information. These weights discount
measurements sampled at the original noise level, rather than changing that
sampling noise.

Geometric sharing averages precisions and information vectors. With
$\bar\mu_{ik,t}=\bar\Pi_{ik,t}^{-1}\bar h_{ik,t}$, the predictive score is
\begin{equation}
 V_{ik,t}=\sigma^2I_3+
    \widetilde H_{i,t}\bar\Pi_{ik,t}^{-1}\widetilde H_{i,t}^{\top},\qquad
 \ell_{ik,t}=\log\Normal(y_{i,t};\widetilde H_{i,t}\bar\mu_{ik,t},V_{ik,t}).
 \label{eq:gaussian-score}
\end{equation}
It accounts for both prediction error and uncertainty. Learning adds the
likelihood's quadratic and linear terms to $(\bar\Pi,\bar h)$; memory then
averages both parameters with their references. Appendix~\ref{app:gaussian}
gives the full updates. For $\omega<1$, accumulated precision is bounded when
observation weights are bounded, so reference beliefs remain influential.

\section{Experimental design}
\label{sec:design}
\subsection{Two explanatory accounts and a common signal}
\label{sec:world}
The shared observables are combustion extent $C$, net gas uptake $G$, and change
in solid mass $M$. The release account introduces an unobserved quantity $P$
representing hypothetical phlogiston release; the oxygen account introduces
unobserved oxygen incorporation $O$. The directed dependencies are shown in
Fig.~\ref{fig:model}. Each coefficient has magnitude $0.8$, with a negative
coefficient only for $P\to M$. Coordinates are standardized, and the hidden
quantities are centred at zero. Each candidate stores
$x_k=(\zeta_k,C,G,M)^\top$, with $\zeta_P=P$ or $\zeta_O=O$, and
$H=[\,0_{3\times1}\ I_3\,]$ selects the observables.

\begin{figure}[!htbp]
 \centering\includegraphics[width=0.92\linewidth]{fig_model.pdf}
 \caption{Reference dependencies for the release (a) and oxygen (b) accounts.
 Dashed nodes are unobserved explanatory quantities. Arrows give regression
 coefficients with the displayed signs and magnitude $0.8$. Both accounts
 connect combustion to mass change through an explanatory substance. Variables
 are standardized schematic quantities.}
 \label{fig:model}
\end{figure}

Reference residual variances are $0.25$ for the explanatory quantity and $0.5$
for each observable. The reference observable means are $(1,1,-1)^\top$ for
release and $(1,1,1)^\top$ for oxygen. References are shared across agents within
each candidate, and initial candidate probabilities are $1/2$. We use $P$ and
$O$ as candidate labels in subscripts as well as names for their hidden quantities.
Appendix~\ref{app:parameters} specifies the reference matrices and intercepts.

The changed-signal experiments use
\begin{equation}
 u_t=\begin{cases}
 (1,1,-1)^\top,&0\leq t<80,\\
 (1,1,1)^\top,&80\leq t<T.
 \end{cases}
 \label{eq:signal}
\end{equation}
Only the observed mass signal changes. The intervention represents a change in
the evidence available to the community. These are repeated noisy measurements
of a prescribed state. The constant-world comparison uses $(1,1,1)^\top$
throughout. A separate experiment with fresh states drawn from the candidate
models is described in Appendix~\ref{app:dependence}.

\subsection{Configurations and measurements}
\label{sec:measurements}
The principal configuration has two candidates, $\omega=0.9$, and
$\alpha_m=0.02$. For an undirected contact graph with adjacency $A$,
$W^x=W^m=W$ has entries
$W_{ij}=(A_{ij}+\delta_{ij})/(\deg(i)+1)$, where $\delta_{ij}$ is one for
$i=j$ and zero otherwise. A complete graph gives $W_{ij}=1/N$.
Single-Gaussian controls use the same conditional update with $K=1$ and references
constructed on the shared observable space. A further control reduces a
neighbour's influence when its mean differs from the receiver's mean; this
social reliability gate is defined in Appendix~\ref{app:gate}.

For a reported agent set $\mathcal I$, the post-round summaries are
\begin{equation}
 \bar q_{\mathcal I,t}(O)=\frac1{|\mathcal I|}\sum_{i\in\mathcal I}q_{iO,t+1},\qquad
 \bar{\mathcal H}_{\mathcal I,t}=\frac1{|\mathcal I|}\sum_{i\in\mathcal I}
       \left[-\sum_kq_{ik,t+1}\log q_{ik,t+1}\right].
 \label{eq:readouts}
\end{equation}
Entropy measures uncertainty over candidates, in nats, with $0\log0=0$. For two
candidates it reaches $\log2$ at equal probabilities. Computing entropy within
each agent before averaging distinguishes individual uncertainty from differences
between confident agents. A crossing time is the first plotted index with
$\bar q_{\mathcal I,t}(O)>0.9$, restricted to $t\geq80$ after a prehistory.
A run without a crossing is recorded as right-censored: observation ended before
the threshold was reached. Individual crossings are recorded separately.

Prediction is assessed on all three observables. Mean squared error (MSE) compares
the probability-weighted candidate mean with the noiseless signal, averaged over
coordinates. Negative log likelihood (NLL) evaluates the full predictive mixture
on four independent new observations per agent per round; lower values indicate
better predictive density. Both are measured after learning and memory, separately
from the score that updates candidate weights. Numerical results average the last
quarter of each run unless stated otherwise. Figure intervals are pointwise
95\% percentile-bootstrap intervals for seed means, using 2,000 resamples of
whole seed trajectories. Paired comparisons resample within-seed differences.
Appendix~\ref{app:protocol} gives the exact estimators and random-seed protocol.

\section{Results}
\label{sec:results}
\subsection{Delayed revision after a change in evidence}
\label{sec:cycle}
Sixteen fully observing agents communicate on a complete graph for 240 rounds
with $\sigma=1$, $\omega=0.9$, and 50 seeds. From equal initial model probabilities,
the community first favours release. After the mass signal changes, support
moves toward oxygen through a period of increased individual uncertainty
(Fig.~\ref{fig:cycle}). All 50 group means cross $0.9$, with median crossing
index $204.5$, or $124.5$ rounds after the change. Of 800 agents, 785 cross
individually. Final mean oxygen probability is $0.9689$.

\begin{figure}[!htbp]
 \centering\includegraphics[width=\linewidth]{fig04_cycle.pdf}
 \caption{Delayed revision in the two-candidate model ($N=16$, 50 seeds).
 Green shows mean oxygen probability on the left axis; grey shows mean individual
 entropy on the right. Narrow ribbons give pointwise 95\% bootstrap intervals
 for seed means. The dotted line marks the first changed observation at $t=80$.
 The amber interval ($t=100$--$188$) marks mean individual entropy above $0.35$
 nats after the change; it is a descriptive threshold.}
 \label{fig:cycle}
\end{figure}

The unchanged-world control has no group or individual crossings after round 80.
With conditional forgetting disabled ($\omega=1$), there are no group crossings
within 240 rounds, although final oxygen probability reaches $0.721$. The
single-Gaussian control changes its numerical predictions. Its last-quarter MSE
is $0.01819$, compared with $0.01009$ for two candidates; the paired difference
is $-0.00810$ with 95\% interval $[-0.00867,-0.00758]$. Corresponding NLLs are
$4.29063$ and $4.27959$. Thus the candidate configuration improves prediction in
this comparison while also providing an explicit measure of theoretical
commitment.

\subsection{Communication under selective observation}
\label{sec:adaptation}
This experiment gives all 16 agents the same fully observed, communicating
prehistory for 80 rounds. At the signal change, eight agents become gas-only
observers, $w=(0,1,0)$, and eight become mass-only observers, $w=(0,0,1)$.
Each seed branches from the identical pre-change state into selective observation
with both communication channels, selective observation in isolation, or independent
full observation. The last condition gives each agent all three channels after
the change, without communication. Runs use $T=1500$, $\sigma=1$, $\omega=0.9$,
and 20 paired seeds.

\begin{figure}[!htbp]
 \centering\includegraphics[width=\linewidth]{fig_adaptation.pdf}
 \caption{Adaptation after a common prehistory (20 seeds). (a,b) Mean oxygen
 probabilities by group for the first 500 of 1,500 rounds. ``Full information''
 denotes independent agents observing all three channels. (c,d) MSE and predictive
 NLL averaged over the final 375 rounds; lower is better. NLL is in nats per
 three-coordinate observation. Ribbons and whiskers give pointwise 95\%
 whole-seed bootstrap intervals.}
 \label{fig:adaptation}
\end{figure}

With communication, all 320 agents cross the individual threshold. In isolation,
all 160 mass observers cross, while none of the 160 gas observers does; the latter
group ends at mean oxygen probability $0.00144$. Communication also improves
full-observable prediction (Table~\ref{tab:adaptation}). The paired NLL difference,
sharing minus isolation, is $-0.6310$ nats, with 95\% interval
$[-0.6332,-0.6290]$. The gas observers lack the changed mass measurement, but
receive its consequences through other agents' beliefs. The independent
full-observation reference crosses sooner. The number and source of its observations differ from those of the communicating
condition, which also draws on neighbours' information.

\begin{table}[!htbp]
 \centering\small
 \caption{Adaptation after the signal change. Crossings refer to group means;
 the median is over observed crossings. Losses average the final 375 rounds.
 A dash denotes no crossing within the run.}
 \label{tab:adaptation}
 \begin{tabularx}{\linewidth}{@{}Yrrrr@{}}
 \toprule
 Condition & Crossings & Median $t$ & MSE & NLL\\
 \midrule
 Selective, sharing & 20/20 & 246 & 0.00979 & 4.30479\\
 Selective, isolated & 0/20 & --- & 0.68184 & 4.93580\\
 Independent, full observation & 20/20 & 217 & 0.03575 & 4.31292\\
 \bottomrule
 \end{tabularx}
\end{table}

\FloatBarrier
\subsection{Network structure and access to evidence}
\label{sec:networks}
The network comparison has 24 agents, 600 rounds, and 50 seeds, with $\sigma=1$
and $\omega=0.9$. After a fully observed 80-round prehistory, four agents retain
full attention and the other twenty use $w=(1,0.01,0.01)$. Figure~\ref{fig:topology}
reports the restricted group. All 50 group means cross in the complete, wheel,
and two-community graphs, with median indices $248$, $261$, and $530$,
respectively. None crosses on the ring within the run; final mean oxygen
probability is $0.735$ and is still increasing.

\begin{figure}[!htbp]
 \centering\includegraphics[width=\linewidth]{fig06_topology.pdf}
 \caption{Network comparison ($N=24$, 50 seeds). Black outlines mark the four
 fully observing agents. Heavy curves and pointwise 95\% bootstrap bands summarize
 the restricted-group seed means; thin curves show its individual agents in seed
 0. ``Cycle'' denotes the ring graph. The dotted line marks $t=80$.
 $\lambda_2$ is the second eigenvalue of the unnormalized graph Laplacian, defined
 in Appendix~\ref{app:networks}. The ring has no group crossing by $t=599$.}
 \label{fig:topology}
\end{figure}

Moving the wheel's fully observing agents from a set including the hub to four
rim nodes delays the median crossing to $447$. Thus source placement matters
within a fixed graph. Across graphs, density and source influence also vary;
the eigenvalue displayed in the figure is a graph descriptor rather than an
isolated explanation of the observed ordering.

On the complete graph, disabling model-weight sharing while retaining conditional
sharing leaves the final restricted-group probability at $0.146$, compared with
$0.772$ when both channels are disabled. Conditional sharing nevertheless gives
better prediction than isolation: NLL is $4.6155$ versus $4.6659$. Full sharing
reduces NLL to $4.3073$. Theory probability and predictive quality therefore need
to be assessed separately. Appendix~\ref{app:controls} reports the additional
communication and observation controls.

\FloatBarrier
\subsection{Agreement on theories and disagreement within them}
\label{sec:camps}
In the constant-world comparison, two groups of eight agents observe only gas
uptake or only mass change for 1,500 rounds ($\sigma=0.5$, 20 seeds).
With two candidates, complete communication, and $\omega=0.3$, both groups
ultimately favour oxygen. Their selected-candidate observable means retain a
last-quarter distance of $0.1083$, and $0.1092$ in a continuation to 5,000 rounds.
This residual difference concerns quantities within the same preferred account.
Appendix~\ref{app:camps} presents the distance and candidate probabilities together.

The plain, gated, and two-candidate configurations use retention values $1$,
$0.95$, and $0.3$ in this comparison. With retention matched at $0.9$, their
last-quarter MSEs are $0.03247$, $0.03704$, and $0.00343$, respectively.
Reference construction still differs between single- and two-candidate models.
The dependence of the distance comparison on memory and reference beliefs makes
it an unsuitable stand-alone measure of the benefit of representing alternatives.

\section{Discussion}
\label{sec:discussion}
Communication changes a candidate's working predictions and transmits support
for it. Under selective observation, these exchanges enable agents who lack the
changed mass measurement to revise their preferred explanation and improve their
predictions.

The cycle gives a computational account of sustained commitment, uncertainty,
and revision \citep{kuhn1962}. With two candidates, entropy is maximal at equal
probabilities. Its rise describes the transition rather than independently
establishing a historical crisis. Linked commitments enter
through reference dependencies, while all quantities follow the same conditional
update.

Agreement on an explanation can coexist with disagreement about its quantities,
and accurate prediction with uncertain theory choice. Appendix~\ref{app:dynamics}
relates these distinctions to the shared and differing components of log odds,
complementing work on individual and collective generative models \citep{waade2025}.

The main experiments compare candidates with different means and dependencies.
The separate fixed-model comparison matches individual observable distributions
and tests information in their dependencies (Appendix~\ref{app:dependence}).
Generating candidates could address unconceived alternatives \citep{stanford2006};
active-inference structural search provides possible approaches
\citep{smith2020structure,neacsu2022}.

Social pooling leaves dependencies between exchanged evidence untracked, and
memory discounts information; this differs from centralized conditioning.
Attention and the signal change are prescribed. Expected free energy could
extend the model to experimental choice \citep{friston2015,millidge2021}.
Results concern finite horizons and the tested configurations: noncrossings
remain censored, and matched retention does not also match reference construction.

\paragraph{Conclusion.}
Within-theory learning and between-theory revision have distinct social
consequences. Communication enables revision under selective observation, with
timing affected by network structure and source placement. Agreement on a theory
can coexist with persistent differences in its estimated quantities. Candidate
probabilities, individual uncertainty, and predictive performance distinguish
these outcomes.
\label{endmain}

\clearpage
\begin{thebibliography}{99}
\bibitem[Bala and Goyal(1998)]{balagoyal1998}
V.~Bala and S.~Goyal.
Learning from neighbours. \emph{Review of Economic Studies}, 65(3):595--621, 1998.

\bibitem[Zollman(2010)]{zollman2010}
K.~J.~S.~Zollman.
The epistemic benefit of transient diversity. \emph{Erkenntnis}, 72(1):17--35, 2010.

\bibitem[Zollman(2013)]{zollman2013}
K.~J.~S.~Zollman.
Network epistemology: communication in epistemic communities.
\emph{Philosophy Compass}, 8(1):15--27, 2013.

\bibitem[DeGroot(1974)]{degroot1974}
M.~H.~DeGroot.
Reaching a consensus. \emph{Journal of the American Statistical Association}, 69(345):118--121, 1974.

\bibitem[Genest and Zidek(1986)]{genest1986}
C.~Genest and J.~V.~Zidek.
Combining probability distributions: a critique and an annotated bibliography.
\emph{Statistical Science}, 1(1):114--135, 1986.

\bibitem[Quine(1951)]{quine1951}
W.~V.~O.~Quine.
Two dogmas of empiricism.
\emph{The Philosophical Review}, 60(1):20--43, 1951.

\bibitem[Lakatos(1970)]{lakatos1970}
I.~Lakatos.
Falsification and the methodology of scientific research programmes.
In I.~Lakatos and A.~Musgrave, editors, \emph{Criticism and the Growth of Knowledge},
pages 91--196. Cambridge University Press, 1970.

\bibitem[Kuhn(1962)]{kuhn1962}
T.~S.~Kuhn.
\emph{The Structure of Scientific Revolutions}. University of Chicago Press, 1962.

\bibitem[Laudan(1977)]{laudan1977}
L.~Laudan.
\emph{Progress and Its Problems: Towards a Theory of Scientific Growth}.
University of California Press, 1977.

\bibitem[Bovens and Hartmann(2003)]{bovens2003}
L.~Bovens and S.~Hartmann.
\emph{Bayesian Epistemology}. Oxford University Press, 2003.

\bibitem[Thagard(1989)]{thagard1989}
P.~Thagard.
Explanatory coherence. \emph{Behavioral and Brain Sciences}, 12(3):435--467, 1989.

\bibitem[Thagard(1990)]{thagard1990}
P.~Thagard.
The conceptual structure of the chemical revolution.
\emph{Philosophy of Science}, 57(2):183--209, 1990.

\bibitem[Kemp and Tenenbaum(2008)]{kemp2008form}
C.~Kemp and J.~B.~Tenenbaum.
The discovery of structural form.
\emph{Proceedings of the National Academy of Sciences}, 105(31):10687--10692, 2008.

\bibitem[Ullman and Tenenbaum(2020)]{ullman2020}
T.~D.~Ullman and J.~B.~Tenenbaum.
Bayesian models of conceptual development: learning as building models of the world.
\emph{Annual Review of Developmental Psychology}, 2:533--558, 2020.

\bibitem[Albarracin et al.(2022)]{albarracin2022}
M.~Albarracin, D.~Demekas, M.~J.~D.~Ramstead, and C.~Heins.
Epistemic communities under active inference. \emph{Entropy}, 24(4):476, 2022.

\bibitem[Friston et al.(2024a)]{friston2023federated}
K.~J.~Friston, T.~Parr, C.~Heins, A.~Constant, D.~Friedman, T.~Isomura, C.~Fields,
T.~Verbelen, M.~Ramstead, J.~Clippinger, and C.~D.~Frith.
Federated inference and belief sharing.
\emph{Neuroscience \& Biobehavioral Reviews}, 156:105500, 2024.

\bibitem[Hyland and Albarracin(2026)]{hyland2025}
D.~Hyland and M.~Albarracin.
On the variational costs of changing our minds.
In \emph{Active Inference: IWAI 2025}, CCIS vol.~2857, pages 3--23.
Springer, 2026. doi:10.1007/978-3-032-16955-6\_1.
Preprint: arXiv:2509.17957, 2025.

\bibitem[\c{C}atal et al.(2024)]{catal2024}
O.~\c{C}atal, T.~Van~de~Maele, R.~J.~Pitliya, M.~Albarracin, C.~Pattisapu, and T.~Verbelen.
Belief sharing: a blessing or a curse. arXiv:2407.02465, 2024.

\bibitem[Waade et al.(2025)]{waade2025}
P.~T.~Waade, C.~L.~Olesen, J.~E.~Laursen, S.~W.~Nehrer, C.~Heins, K.~Friston, and C.~Mathys.
As one and many: relating individual and emergent group-level generative models in active inference.
\emph{Entropy}, 27(2):143, 2025.

\bibitem[Jadbabaie et al.(2012)]{jadbabaie2012}
A.~Jadbabaie, P.~Molavi, A.~Sandroni, and A.~Tahbaz-Salehi.
Non-Bayesian social learning. \emph{Games and Economic Behavior}, 76(1):210--225, 2012.

\bibitem[Molavi et al.(2018)]{molavi2018}
P.~Molavi, A.~Tahbaz-Salehi, and A.~Jadbabaie.
A theory of non-Bayesian social learning. \emph{Econometrica}, 86(2):445--490, 2018.

\bibitem[Nedi\'c et al.(2017)]{nedic2017}
A.~Nedi\'c, A.~Olshevsky, and C.~A.~Uribe.
Fast convergence rates for distributed non-Bayesian learning.
\emph{IEEE Transactions on Automatic Control}, 62(11):5538--5553, 2017.

\bibitem[Lalitha et al.(2018)]{lalitha2018}
A.~Lalitha, T.~Javidi, and A.~D.~Sarwate.
Social learning and distributed hypothesis testing.
\emph{IEEE Transactions on Information Theory}, 64(9):6161--6179, 2018.

\bibitem[Hare et al.(2019)]{hare2019gaussian}
J.~Z.~Hare, C.~Uribe, L.~Kaplan, and A.~Jadbabaie.
Non-Bayesian social learning with Gaussian uncertain models.
arXiv:1910.11251, 2019.

\bibitem[Friston et al.(2017b)]{friston2017process}
K.~Friston, T.~FitzGerald, F.~Rigoli, P.~Schwartenbeck, and G.~Pezzulo.
Active inference: a process theory. \emph{Neural Computation}, 29(1):1--49, 2017.

\bibitem[Da Costa et al.(2020)]{dacosta2020}
L.~Da Costa, T.~Parr, N.~Sajid, S.~Veselic, V.~Neacsu, and K.~Friston.
Active inference on discrete state-spaces: a synthesis.
\emph{Journal of Mathematical Psychology}, 99:102447, 2020.

\bibitem[Blumenthal and Ladyman(2017)]{blumenthal2017}
G.~Blumenthal and J.~Ladyman.
The development of problems within the phlogiston theories, 1766--1791.
\emph{Foundations of Chemistry}, 19(3):241--280, 2017.

\bibitem[Holmes(2000)]{holmes2000}
F.~L.~Holmes.
The revolution in chemistry and physics: overthrow of a reigning paradigm or competition between
contemporary research programs? \emph{Isis}, 91(4):735--753, 2000.

\bibitem[Boantza and Gal(2011)]{boantza2011}
V.~D.~Boantza and O.~Gal.
The `absolute existence' of phlogiston: the losing party's point of view.
\emph{British Journal for the History of Science}, 44(3):317--342, 2011.

\bibitem[Stanford(2006)]{stanford2006}
P.~K.~Stanford.
\emph{Exceeding Our Grasp: Science, History, and the Problem of Unconceived Alternatives}.
Oxford University Press, 2006.

\bibitem[Smith et al.(2020)]{smith2020structure}
R.~Smith, P.~Schwartenbeck, T.~Parr, and K.~J.~Friston.
An active inference approach to modeling structure learning: concept learning as an example case.
\emph{Frontiers in Computational Neuroscience}, 14:41, 2020.

\bibitem[Neacsu et al.(2022)]{neacsu2022}
V.~Neacsu, M.~B.~Mirza, R.~A.~Adams, and K.~J.~Friston.
Structure learning enhances concept formation in synthetic active inference agents.
\emph{PLOS ONE}, 17(11):e0277199, 2022.

\bibitem[Friston et al.(2015)]{friston2015}
K.~Friston, F.~Rigoli, D.~Ognibene, C.~Mathys, T.~Fitzgerald, and G.~Pezzulo.
Active inference and epistemic value. \emph{Cognitive Neuroscience}, 6(4):187--214, 2015.

\bibitem[Millidge et al.(2021)]{millidge2021}
B.~Millidge, A.~Tschantz, and C.~L.~Buckley.
Whence the expected free energy? \emph{Neural Computation}, 33(2):447--482, 2021.
\end{thebibliography}
\clearpage
\appendix
\renewcommand{\theHsection}{appendix.\Alph{section}}
\renewcommand{\theHsubsection}{\theHsection.\arabic{subsection}}
\section{Distributional updates}
\label{app:distribution}
\subsection{Notation and normalization}
Table~\ref{tab:notation} collects the notation. State $t$ enters round $t$;
barred and $+$ quantities refer to that round, and the stored result is state
$t+1$. Only the logarithmic camp-distance axis uses completed updates $t+1$.

\begin{table}[!htbp]
 \centering\small
 \caption{Notation. Matrix and vector dimensions refer to the experimental
 implementation; the distributional formulation permits candidate-specific spaces.}
 \label{tab:notation}
 \begin{tabularx}{\linewidth}{@{}p{3.4cm}Y@{}}
 \toprule
 Symbol & Meaning\\
 \midrule
 $i,j$; $k,k'$; $t,T$ & Agents; candidates; round and run length.\\
 $x_k$, $Q_{i,t}$ & Candidate state; joint candidate--state density.\\
 $\rho_{ik,t}$, $\rho^0_{ik}$ & Working and fixed reference distributions.\\
 $q_{ik,t}$, $L_{ik,t}$ & Candidate probability and log weight.\\
 $\bar{(\cdot)}$, $(\cdot)^+$ & After sharing; after learning, before memory.\\
 $W^x_t$, $W^m_t$; $\alpha_m,\omega$ & Sharing matrices; model-sharing rate and retention.\\
 $o_{i,t}$, $y_{i,t}$, $u_t$; $\sigma$ & Raw and weighted observations, signal; noise standard deviation.\\
 $w_{i,t}$, $S_{i,t}$, $H$ & Channel weights, their diagonal square root, observation map.\\
 $\mu_{ik,t},\Sigma_{ik,t}$; $\Pi_{ik,t},h_{ik,t}$ & Mean and covariance; inverse covariance and information vector.\\
 $J_{i,t}$, $j_{i,t}$; $V_{ik,t}$ & Likelihood increments; predictive covariance.\\
 $B_k$, $b_k$, $R_k$ & Directed coefficients, intercepts, residual covariance.\\
 $\mathcal I$, $\hat u_{i,t}$, $D_t$ & Reported group, mixture-mean prediction, camp distance.\\
 \bottomrule
 \end{tabularx}
\end{table}

For each candidate, densities use a common base measure across agents. The
conditional sharing normalizer is
\begin{equation}
 Z^x_{ik,t}=\int\prod_j\rho_{jk,t}(x_k)^{W^x_{ij,t}}\,dx_k.
 \label{eq:normalizer}
\end{equation}
Generalized H\"older's inequality gives $Z^x_{ik,t}\leq1$ for proper densities
and weights summing to one. Common positive support gives $Z^x_{ik,t}>0$ and
also permits memory interpolation. Nonsingular Gaussians have full support.
The predictive density in Eq.~\eqref{eq:score} must be finite and positive.

Conditional sharing minimizes
$\sum_jW^x_{ij,t}D_{\mathrm{KL}}(\varrho\|\rho_{jk,t})$ over normalized
$\varrho$. Substitution of Eq.~\eqref{eq:normalizer} makes this objective
$D_{\mathrm{KL}}(\varrho\|\bar\rho_{ik,t})-\log Z^x_{ik,t}$.
Memory has the same interpretation with weights $1-\omega$ and $\omega$.
With undamped equal sharing weights for both components, pooling whole joint
distributions instead would give
$q^{\mathrm{joint}}_{ik}\propto Z^x_{ik,t}\prod_jq_{jk,t}^{W_{ij,t}}$.
The implemented, separately normalized pools in Eq.~\eqref{eq:sharing} express
conditional agreement and model support through distinct operations.

\subsection{Joint local inference}
Let $\bar q_{ik,t}\propto\exp\bar L_{ik,t}$ be normalized over candidates.
For trial probabilities $\pi_k$ and conditional densities $\varrho_k$, the
local objective is
\begin{equation}
 \mathcal F_{i,t}[\pi,\{\varrho_k\}]
   =\KL{\pi}{\bar q_{i,t}}+\sum_k\pi_k\mathcal F_{ik,t}[\varrho_k].
 \label{eq:joint-free-energy}
\end{equation}
The minimizers $\varrho_k=\rho^+_{ik,t}$ and
$\pi_k\propto\bar q_{ik,t}\exp\ell_{ik,t}$ recover Eq.~\eqref{eq:learning}.
Restricting the conditional family adds approximation error to each optimized
free energy. Those errors may differ between candidates; the Gaussian experiments
use exact predictive evidence.

\section{Gaussian implementation and reference models}
\label{app:gaussian}
\subsection{Complete update and precision bound}
A Gaussian log density has quadratic term $-x^\top\Pi x/2$ and linear term
$h^\top x$. Collecting these terms under sharing, learning, and memory gives
\begin{align}
 \bar\Pi_{ik,t}&=\sum_j W^x_{ij,t}\Pi_{jk,t},&
 \bar h_{ik,t}&=\sum_j W^x_{ij,t}h_{jk,t},\nonumber\\
 J_{i,t}&=\sigma^{-2}H^\top\diag(w_{i,t})H,&
 j_{i,t}&=\sigma^{-2}H^\top\diag(w_{i,t})o_{i,t},\nonumber\\
 \Pi^+_{ik,t}&=\bar\Pi_{ik,t}+J_{i,t},&
 h^+_{ik,t}&=\bar h_{ik,t}+j_{i,t},\nonumber\\
 \Pi_{ik,t+1}&=(1-\omega)\Pi^0_{ik}+\omega\Pi^+_{ik,t},&
 h_{ik,t+1}&=(1-\omega)h^0_{ik}+\omega h^+_{ik,t}.
 \label{eq:natural-update}
\end{align}
Both candidates receive the same likelihood increments. Their pre-update means
and covariances, and hence their predictive scores, can differ. Memory discounts
the current increment as well as older information because it follows learning.


With $e_{ik,t}=y_{i,t}-\widetilde H_{i,t}\bar\mu_{ik,t}$, the experimental
three-channel score expands to
\begin{equation}
 \ell_{ik,t}=-\tfrac12\left[3\log(2\pi)+\log\det V_{ik,t}
                          +e_{ik,t}^\top V_{ik,t}^{-1}e_{ik,t}\right].
 \label{eq:score-expanded}
\end{equation}
Linear solves and triangular factorizations evaluate the score. Proper predictive
densities follow from positive reference precisions and positive noise variance.

Fractional attention has a precise interpretation. For positive weights,
\begin{equation}
 \Normal(So;SHx,\sigma^2I_3)
  =\frac{\Normal(o;Hx,\sigma^2\diag(w)^{-1})}{\det S}.
 \label{eq:attention-jacobian}
\end{equation}
The denominator is candidate-independent, so it cancels in model odds.
A zero-weight row has zero transformed observation and observation operator,
contributes the same score constant to every candidate, and adds no information.
The transformed sampling noise has covariance $\sigma^2\diag(w)$, whereas the
working likelihood uses $\sigma^2I_3$. Thus inference deliberately discounts
the sampled observations.

For the matrix operator norm, set $p_t=\max_{i,k}\|\Pi_{ik,t}\|$,
$p_{\mathrm{ref}}=\max_{i,k}\|\Pi^0_{ik}\|$, and suppose
$\|J_{i,t}\|\leq J_{\max}$. Convexity gives, for $0\leq\omega<1$,
\begin{equation}
 p_t\leq\omega^t p_0+(1-\omega^t)
       \left(p_{\mathrm{ref}}+\frac{\omega J_{\max}}{1-\omega}\right).
 \label{eq:precision-bound}
\end{equation}
This follows by iterating
$p_{t+1}\leq(1-\omega)p_{\mathrm{ref}}+\omega p_t+\omega J_{\max}$.
An analogous deterministic bound on $h$ would require bounded $j$; Gaussian
sampling noise is not uniformly bounded. The precision bound controls information
accumulation, not the sign of candidate log odds.

\subsection{Numerical reference specification}
\label{app:parameters}
A reference directed model obeys $x_k=B_kx_k+b_k+\varepsilon_k$, with independent
Gaussian residuals $\varepsilon_k\sim\Normal(0,R_k)$. Entry $(B_k)_{ab}$ is the
coefficient from coordinate $b$ to coordinate $a$. For $F_k=I_4-B_k$,
\begin{equation}
 \Pi^0_k=F_k^\top R_k^{-1}F_k,\quad h^0_k=F_k^\top R_k^{-1}b_k,\quad
 \mu^0_k=F_k^{-1}b_k,\quad \Sigma^0_k=F_k^{-1}R_kF_k^{-\top}.
 \label{eq:compile}
\end{equation}
The directed graphs are acyclic, so $F_k$ is invertible. Residual variances are
positive, giving positive-definite $\Pi^0_k$. Candidate references in the
experiments are independent of agent index $i$.

In state order $(\zeta_k,C,G,M)$ the exact coefficients are
\begin{equation}
 B_P=\begin{pmatrix}0&.8&0&0\\0&0&0&0\\0&.8&0&0\\-.8&0&0&0\end{pmatrix},\qquad
 B_O=\begin{pmatrix}0&.8&0&0\\0&0&0&0\\.8&0&0&0\\.8&0&0&0\end{pmatrix}.
 \label{eq:B}
\end{equation}
Both have $R_k=\diag(.25,.5,.5,.5)$, and
\begin{equation}
 \begin{aligned}
 \mu^0_P&=(0,1,1,-1)^\top,& b_P&=(-.8,1,.2,-1)^\top,\\
 \mu^0_O&=(0,1,1,1)^\top,& b_O&=(-.8,1,1,1)^\top.
 \end{aligned}
 \label{eq:means-intercepts}
\end{equation}
Intercepts satisfy $b_k=F_k\mu^0_k$; they differ from joint means.
The reference observable covariances are
\begin{equation}
 \Sigma_P^{u,0}=\begin{pmatrix}.5&.4&-.32\\.4&.82&-.256\\-.32&-.256&.8648\end{pmatrix},\quad
 \Sigma_O^{u,0}=\begin{pmatrix}.5&.32&.32\\.32&.8648&.3648\\.32&.3648&.8648\end{pmatrix}.
 \label{eq:observable-covariances}
\end{equation}
The superscript $u$ denotes the observable marginal. Conditional on the
underlying $C$, the reference covariance between $G$ and $M$ is zero for release
and $0.16$ for oxygen. Learning updates full Gaussian parameters; the directed
coefficients specify their references.

\subsection{Single-Gaussian controls}
\label{app:baselines}
For each candidate, marginalizing to the observables gives
$\mu_k^{u,0}=H\mu_k^0$ and $\Sigma_k^{u,0}=H\Sigma_k^0H^\top$.
Write $\Pi_k^{u,0}=(\Sigma_k^{u,0})^{-1}$ and
$h_k^{u,0}=\Pi_k^{u,0}\mu_k^{u,0}$. The shared single-Gaussian reference is
\begin{equation}
 \Pi_{\mathrm{mid}}^u=\tfrac12(\Pi_P^{u,0}+\Pi_O^{u,0}),\qquad
 h_{\mathrm{mid}}^u=\tfrac12(h_P^{u,0}+h_O^{u,0}).
 \label{eq:midpoint}
\end{equation}
The cycle control starts at this reference. In the camp controls, group A starts
at the release observable marginal and group B at the oxygen marginal. Their
references remain group-specific or use Eq.~\eqref{eq:midpoint}, as listed in
Table~\ref{tab:camp-controls}. An independent unobserved coordinate of variance
$0.25$ keeps arrays four-dimensional and is excluded from comparisons.

\begin{figure}[!htbp]
 \centering\includegraphics[width=\linewidth]{fig02_precision_pool.pdf}
 \caption{Construction of the shared single-Gaussian reference. Links show
 off-diagonal entries of the observable marginal precision matrices and their
 equal-weight average, rounded to three decimals. Solid blue and dashed orange
 links are negative and positive entries, respectively. Diagonals and information
 vectors are included in the calculation but omitted from the drawing.}
 \label{fig:precision}
\end{figure}

An observable marginal precision is the inverse of the observable covariance,
not the principal block of the full precision. In a partition into observable
$u$ and hidden $\zeta$, it equals
$\Pi_{uu}-\Pi_{u\zeta}\Pi_{\zeta\zeta}^{-1}\Pi_{\zeta u}$; the corresponding
information vector is $h_u-\Pi_{u\zeta}\Pi_{\zeta\zeta}^{-1}h_\zeta$.
The midpoint therefore combines comparable observable beliefs without equating
$P$ with $O$.

Precision links in Fig.~\ref{fig:precision} are undirected conditional
associations. In a Gaussian, a zero off-diagonal precision entry indicates
conditional independence given the other coordinates. Directed reference
coefficients have a different meaning. Pooling preserves an entry that is zero
in every input, while opposite-signed entries can cancel. Here $H$ selects
coordinates, making $J_{i,t}$ diagonal. A common off-diagonal zero in both the
initial state and the references remains zero under the implemented update.

\subsection{Social reliability gate}
\label{app:gate}
For $K=1$, write $\mu_{i,t}$ and $\Pi_{i,t}$ for the conditional parameters.
Let $\mathcal V=\{2,3,4\}$ index the observable coordinates in the mathematical
state order. The gate uses the entering state:
\begin{equation}
 \begin{gathered}
 d_{ij,t}^2=\sum_{v\in\mathcal V}
 \frac{(\mu_{iv,t}-\mu_{jv,t})^2}
 {1/(\Pi_{i,t})_{vv}+1/(\Pi_{j,t})_{vv}},\qquad
 \gamma_{ij,t}=\frac{\nu_s+1}{\nu_s+d_{ij,t}^2},\\
 W^x_{ij,t}=\frac{(A_{ij}+\delta_{ij})\gamma_{ij,t}}
 {\sum_a(A_{ia}+\delta_{ia})\gamma_{ia,t}}.
 \end{gathered}
 \label{eq:gate}
\end{equation}
Here $\nu_s=0.01$. A larger disagreement reduces relative influence. The
reciprocal $1/(\Pi_i)_{vv}$ is the Gaussian variance conditional on the other
coordinates, whereas $(\Pi_i^{-1})_{vv}$ is its marginal variance. The gate
uses the former and omits cross-coordinate covariance. Its denominator includes
$10^{-12}$ in the implementation for numerical protection. This control acts on
social messages; $w$ separately weights incoming observations. Two-candidate
runs use ungated sharing.

\section{Measurements and experimental protocol}
\label{app:protocol}
\subsection{Predictive evaluation and candidate readouts}
Define the post-round observable candidate mean
$\mu^u_{ik,t+1}=H\mu_{ik,t+1}$ and covariance
$\Sigma^u_{ik,t+1}=H\Pi_{ik,t+1}^{-1}H^\top$. The observable prediction and
selected-candidate estimate are
\begin{equation}
 \hat u_{i,t}=\sum_kq_{ik,t+1}\mu^u_{ik,t+1},\qquad
 \hat u^{\mathrm{sel}}_{i,t}=\mu^u_{i\hat k_{i,t},t+1},\qquad
 \hat k_{i,t}=\arg\max_kL_{ik,t+1}.
 \label{eq:observable-readouts}
\end{equation}
Ties select the first candidate in the stored order. The mean squared error is
$\mathrm{MSE}_{i,t}=\|\hat u_{i,t}-u_t\|_2^2/3$. The predictive density is the
mixture
\begin{equation}
 p^{\mathrm{eval}}_{i,t}(o)=\sum_kq_{ik,t+1}
     \Normal(o;\mu^u_{ik,t+1},\Sigma^u_{ik,t+1}+\sigma^2I_3).
 \label{eq:evaluation-mixture}
\end{equation}
Four independent draws $o^{\mathrm{eval}}_{i,t,b}$, $b=1,\ldots,4$, follow
$\Normal(u_t,\sigma^2I_3)$ and are reserved for evaluation. Their loss is
\begin{equation}
 \mathrm{NLL}_{i,t}=-\tfrac14\sum_{b=1}^4
          \log p^{\mathrm{eval}}_{i,t}(o^{\mathrm{eval}}_{i,t,b}).
 \label{eq:nll}
\end{equation}
All channels use the full measurement variance. The logarithm is taken after
mixing the candidate densities.

The distance between groups A and B is
\begin{equation}
 D_t=\left\|\frac1{|\mathrm A|}\sum_{i\in\mathrm A}\hat u^{\mathrm{sel}}_{i,t}
       -\frac1{|\mathrm B|}\sum_{i\in\mathrm B}\hat u^{\mathrm{sel}}_{i,t}
       \right\|_2.
 \label{eq:camp-distance}
\end{equation}
Only shared observables enter this distance. Opposing group preferences are
assessed separately by the simultaneous conditions
$\bar q_{\mathrm A,t}(O)<0.1$ and $\bar q_{\mathrm B,t}(O)>0.9$.
Group and individual threshold crossings use strict inequality and the first
qualifying post-update sample; later recrossings are possible.

For any reported set $\mathcal I$, write $q_i=(q_{ik,t+1})_{k=1}^K$ and
$\bar q=|\mathcal I|^{-1}\sum_{i\in\mathcal I}q_i$ temporarily. Entropy satisfies
\begin{equation}
 \mathcal H(\bar q)-\frac1{|\mathcal I|}\sum_{i\in\mathcal I}\mathcal H(q_i)
     =\frac1{|\mathcal I|}\sum_{i\in\mathcal I}\KL{q_i}{\bar q}\geq0.
 \label{eq:entropy-decomposition}
\end{equation}
Thus population entropy combines individual uncertainty with between-agent
disagreement. Equally sized groups that are certain of opposite candidates
have individual entropy zero and population entropy $\log2$.

\subsection{Design, pairing, and uncertainty}
\label{app:design}
The frozen assessment contains four cycle conditions, nine network conditions,
eleven constant-world conditions, two continuations, and three adaptation
conditions. All two-candidate conditions start from their references with equal
candidate probabilities. Model-score decay is zero throughout. The four cycle
conditions are the changed signal, unchanged signal, no conditional forgetting,
and the single-Gaussian midpoint control. Table~\ref{tab:settings} lists the
common settings; interventions change the stated settings only.

\begin{table}[!htbp]
 \centering\small
 \caption{Common settings by experiment. The world changes at round 80 except
 in the constant-world experiment and unchanged-world control.}
 \label{tab:settings}
 \begin{tabular}{@{}lrrrrr@{}}
 \toprule
 Experiment & $N$ & $T$ & $\sigma$ & $\omega$ & Seeds\\
 \midrule
 Cycle &16&240&1&0.9&$0$--$49$\\
 Networks &24&600&1&0.9&$0$--$49$\\
 Adaptation &16&1500&1&0.9&$0$--$19$\\
 Constant world &16&1500&0.5&Varies&$0$--$19$\\
 \bottomrule
 \end{tabular}
\end{table}

Cycle, constant-world, and adaptation experiments use complete graphs.
All adaptation branches share state 80, after processing observations
$0$--$79$. Complete-graph network controls similarly share their complete-graph
prehistory. For simulation seed $s$, the second-phase observation generator
restarts with seed $s+1000$, identically across paired branches. Evaluation
noise uses a separate random key and does not advance the learning generator.

For a run of length $T$, loss summaries average over the reported agents and
indices $t=\lfloor3T/4\rfloor,\ldots,T-1$, within each seed, and then over seeds.
The group is all agents for cycle, adaptation, and constant-world conditions;
network summaries use the twenty restricted agents. A paired loss contrast
subtracts these seed-level averages before resampling. The 95\% intervals use
the 2.5th and 97.5th percentiles of 2,000 resampled seed means, with bootstrap
seed 20260919. Figure bands resample complete seed trajectories, maintaining
within-run dependence over time. They are pointwise intervals for the mean,
not simultaneous bands or the variation of individual agents.

\subsection{Exact network specification}
\label{app:networks}
Code numbers agents from zero; the mathematical agent index $i$ corresponds to
code index $i-1$. The ring joins consecutive code nodes $0,\ldots,23$ cyclically.
The wheel has hub 0 connected to a ring on nodes $1,\ldots,23$. The community
graph consists of complete blocks $\{0,1,2,3\}$ and $\{4,\ldots,23\}$, joined by
seven edges:
$(0,7)$, $(1,7)$, $(1,19)$, $(2,8)$, $(2,10)$, $(2,22)$, and $(3,23)$.
The four graphs have 24, 46, 203, and 276 undirected edges, respectively.
The archived adjacencies are fixed across simulation seeds.

Nodes $0$--$3$ remain fully observing in the main comparison. The wheel-rim
condition instead assigns those roles to nodes $1$--$4$ by a permutation applied
before simulation. All other nodes use $(1,0.01,0.01)$ after round 80.
For the unnormalized graph Laplacian
$\mathscr L=\diag(\deg(1),\ldots,\deg(N))-A$, the second smallest eigenvalue
$\lambda_2$ appears in Fig.~\ref{fig:topology}. The actual sharing operator is the
row-normalized $W$, whose stationary weights on these undirected graphs satisfy
$v_i=(\deg(i)+1)/\sum_j(\deg(j)+1)$. Source influence therefore varies with
source placement as well as adjacency.

\subsection{Implementation record}
The frozen assessment \path{assessment-2026-09-19} uses source revision
\texttt{a869664}, Python 3.12.9 and JAX/JAXlib 0.9.0 with 64-bit arithmetic.
In \path{paper/combustion_experiments/}, \path{world.py} specifies the models,
\path{run.py} the conditions and evaluation, and \path{structural_check.py}
the separate classifier. The engine's \path{rival.py} implements the update
sequence in Sect.~\ref{sec:update}.
Code state coordinates $(0,1,2,3)$ correspond to mathematical coordinates
$(1,2,3,4)$ in order $(\zeta_k,C,G,M)$; observation coordinates $(0,1,2)$
correspond to $(C,G,M)$.

\section{Additional comparisons}
\label{app:controls}
\subsection{Theory choice and observable distance}
\label{app:camps}
Table~\ref{tab:camp-controls} gives all eleven constant-world conditions. Every
single-Gaussian condition starts from the group-specific observable priors in
Appendix~\ref{app:baselines}, including the conditions with a shared reference.
The two-candidate conditions give both references to every agent. A ``group''
reference means release for gas-observing group A and oxygen for mass-observing
group B. Model-weight sharing has no effect when $K=1$.

\begin{table}[!htbp]
 \centering\small
 \caption{Constant-world conditions: retention, communication, and reference.
 All use 20 seeds, $N=16$, $T=1500$, and $\sigma=0.5$. ``Midpoint'' is
 Eq.~\eqref{eq:midpoint}; ``candidate'' retains each candidate's own reference.}
 \label{tab:camp-controls}
 \begin{tabularx}{\linewidth}{@{}Ycccll@{}}
 \toprule
 Condition & $K$ & $\omega$ & $\alpha_m$ & Conditional sharing & Reference\\
 \midrule
 Plain &1&1&---&Fixed $W$&Group\\
 Gated &1&.95&---&Gated&Group\\
 Two candidates &2&.3&.02&Fixed $W$&Candidate\\
 Plain, matched memory &1&.9&---&Fixed $W$&Group\\
 Gated, matched memory &1&.9&---&Gated&Group\\
 Two candidates, matched memory &2&.9&.02&Fixed $W$&Candidate\\
 Plain, shared reference &1&.9&---&Fixed $W$&Midpoint\\
 Gated, shared reference &1&.9&---&Gated&Midpoint\\
 Two candidates, no forgetting &2&1&.02&Fixed $W$&Candidate\\
 Model-weight sharing off &2&.3&0&Fixed $W$&Candidate\\
 Both channels off &2&.3&0&Identity&Candidate\\
 \bottomrule
 \end{tabularx}
\end{table}

The selected-candidate distance and probability readouts are shown in
Fig.~\ref{fig:camps-distance}. With both channels
and $\omega=0.3$, no seed satisfies the simultaneous opposing-group criterion.
With model-weight sharing disabled but conditional sharing retained, all 20
group means still cross, with median index 109. In isolation, the final gas-group
probability is $0.000655$ and the mass-group probability is approximately one.
Its MSE is $0.68525$, compared with $0.02157$ under full sharing.

\begin{figure}[!htbp]
 \centering
 \includegraphics[width=0.9\linewidth]{fig03_camp_distance.pdf}\\[4pt]
 \includegraphics[width=0.9\linewidth]{fig_camp_choices.pdf}
 \caption{Constant-world comparison (20 seeds). Top: selected-candidate
 observable distance $D_t$, with retention rates in the legend and completed
 updates $t+1$ on the horizontal axis. References also differ across configurations.
 Bottom: candidate probabilities for the same two-candidate, $\omega=0.3$ runs,
 over the first 150 rounds (a) and all 1,500 rounds (b). Both groups favour oxygen.
 Bands give pointwise 95\% whole-seed bootstrap intervals.}
 \label{fig:camps-distance}\label{fig:camps-choice}
\end{figure}

At the original retention rates, two candidates improve point prediction but
have higher NLL (Table~\ref{tab:camp-losses}). With matched retention, they
improve both losses. Without conditional forgetting, final oxygen probability
is $0.8526$ and MSE is $0.00178$: accurate prediction can coexist with less
than decisive candidate support.

\begin{table}[!htbp]
 \centering\small
 \caption{Last-quarter predictive losses in the constant-world experiment.
 Means average the 20 seeds; lower values are better.}
 \label{tab:camp-losses}
 \begin{tabular}{@{}lrrr@{}}
 \toprule
 Configuration & $\omega$ & MSE & NLL\\
 \midrule
 Plain &1&0.02628&2.36789\\
 Gated &.95&0.02890&2.38745\\
 Two candidates &.3&0.02157&2.81750\\
 Plain, matched memory &.9&0.03247&2.41589\\
 Gated, matched memory &.9&0.03704&2.44333\\
 Two candidates, matched memory &.9&0.00343&2.35373\\
 Plain, shared reference &.9&0.03237&2.41524\\
 Gated, shared reference &.9&0.03492&2.42861\\
 Two candidates, no forgetting &1&0.00178&2.33157\\
 \bottomrule
 \end{tabular}
\end{table}

The continuation retains $\omega=0.3$ and extends stored state 1500 by 3,500
rounds, with seed $s+3000$. Its distance $0.1092$ averages the last 875 rounds
of that continuation. A second, 400-round continuation sets all attention to
zero and $\omega=1$, retaining both sharing channels, with seed $s+2000$.
The candidate score differences then vanish and the observed log-odds gap
contracts as derived in Appendix~\ref{app:dynamics}; the maximum numerical
error is $2.04\times10^{-12}$.

\subsection{Network controls}
\label{app:network-controls}
Table~\ref{tab:network-controls} reports the nine network conditions. The group
is always the twenty agents designated as restricted, including the control in
which all agents attend fully. There are 760 noncrossing restricted agents on
the ring and 613 under complete isolation, out of 1,000 agent trajectories in
each condition. These individual counts differ from group crossing counts.

\begin{table}[!htbp]
 \centering\small
 \caption{Network adoption by round 599, over 50 seeds. Medians use observed
 group crossings. A dash denotes no observed crossing. Losses are reported
 in the text separately from adoption.}
 \label{tab:network-controls}
 \begin{tabularx}{\linewidth}{@{}Yrrr@{}}
 \toprule
 Condition & Crossings & Median $t$ & Final $\bar q(O)$\\
 \midrule
 Complete &50/50&248&$\simeq1$\\
 Wheel, hub among full observers &50/50&261&$\simeq1$\\
 Two communities &50/50&530&.9774\\
 Ring &0/50&---&.7349\\
 Wheel, four rim observers &50/50&447&.9993\\
 Complete, model-weight sharing off &0/50&---&.1458\\
 Complete, both channels off &0/50&---&.7723\\
 Complete, restricted weights $(1,0,0)$ &50/50&255&$\simeq1$\\
 Complete, all weights $(1,1,1)$ &50/50&202&$\simeq1$\\
 \bottomrule
 \end{tabularx}
\end{table}

In the zero-direct-gas/mass condition, restricted agents still observe combustion
extent and receive both social channels. Full sharing improves predictive NLL
relative to conditional-only sharing by $0.3082$ nats, with paired 95\% interval
$[0.3033,0.3130]$. The conditional-only condition improves point prediction over
isolation (MSE $0.28723$ versus $0.32576$), even though it produces lower oxygen
probability. These measures describe different consequences of communication.

\section{Log-odds dynamics}
\label{app:dynamics}
For candidates $O$ and $P$, define
$z_{i,t}=L_{iO,t}-L_{iP,t}$ and
$s_{i,t}=\ell_{iO,t}-\ell_{iP,t}$. With vectors $z_t,s_t\in\R^N$,
Eqs.~\eqref{eq:sharing}--\eqref{eq:memory} give
\begin{equation}
 z_{t+1}=P_tz_t+s_t,\qquad P_t=(1-\alpha_m)I_N+\alpha_m W^m_t.
 \label{eq:log-odds}
\end{equation}
The score vector $s_t$ depends on the evolving conditional beliefs and observations.
For a fixed connected undirected graph and $\alpha_m>0$, $P$ has a unique
stationary probability vector $v$, satisfying $v^\top P=v^\top$ and
$v^\top\mathbf1=1$. Positive self-weights ensure convergence of repeated averaging.
Define $c_t=v^\top z_t$ and $d_t=(I_N-\mathbf1v^\top)z_t$, where $\mathbf1$ is
the all-ones vector. Then
\begin{equation}
 c_{t+1}=c_t+v^\top s_t,\qquad
 d_{t+1}=Pd_t+(I_N-\mathbf1v^\top)s_t.
 \label{eq:log-odds-components}
\end{equation}
Communication removes old differences while unequal predictive scores add new
ones. For a constant score vector $s$, the differing component converges to
\begin{equation}
 d_*=(I_N-P+\mathbf1v^\top)^{-1}(I_N-\mathbf1v^\top)s,
 \label{eq:stationary-gap}
\end{equation}
while $c_t$ changes by $v^\top s$ each round. This separates a stationary
difference from stationary candidate probabilities.

For equally sized groups on the complete graph, let
$g_t=\langle z_{i,t}\rangle_{\mathrm A}-\langle z_{i,t}\rangle_{\mathrm B}$ and
$s_{\mathrm A,t},s_{\mathrm B,t}$ be their average score differences. The exact
recurrence is
\begin{equation}
 g_{t+1}=(1-\alpha_m)g_t+s_{\mathrm A,t}-s_{\mathrm B,t}.
 \label{eq:camp-odds}
\end{equation}
For constant sources it has fixed difference
$(s_{\mathrm A}-s_{\mathrm B})/\alpha_m$, for $0<\alpha_m\leq1$.
Both group log odds can nevertheless become positive when the shared component
increases. In the $\omega=0.3$ constant-world runs, the last-quarter difference
is approximately $-45.9$ and the mean score increment is $0.473$ nats per round;
both groups favour oxygen. Their differing quantity estimates and separated
log odds do not imply opposite candidate preferences.

When all observation weights are zero, the candidate score difference is exactly
zero. Starting from state 1500, the continuation therefore has
$g_{1500+n}=(1-0.02)^n g_{1500}$ for $n=1,\ldots,400$.
This test isolates the averaging recurrence from continuing observational input.

\section{Discrimination from cross-variable dependence}
\label{app:dependence}
The main candidate references differ in means and covariance. At measurement-noise
standard deviation $\sigma=1$, the mean-displacement term contributes about
90.7\% and 90.9\% of their two directional initial predictive Gaussian
Kullback--Leibler divergences. These proportions describe initial predictions,
not every score after conditional learning.

A separate experiment tests information in the dependence structure after
matching all individual observable distributions. Both reference means are set
to zero, and the release model's gas residual variance is changed from $0.5$ to
$0.5448$. All other residuals and directed coefficients remain unchanged. The matrices
$R_k$ and $\Sigma_k^0$ in this appendix use that adjustment. Both observable
variance vectors are $(0.5,0.8648,0.8648)$.
Each sample draws a fresh state from generating candidate $k_\star$:
\begin{equation}
 x^{(n)}_{k_\star}=(I_4-B_{k_\star})^{-1}\varepsilon^{(n)},\qquad
 o^{(n)}=Hx^{(n)}_{k_\star}+\sigma\xi^{(n)},
 \label{eq:fresh-samples}
\end{equation}
where $\varepsilon^{(n)}\sim\Normal(0,R_{k_\star})$ and
$\xi^{(n)}\sim\Normal(0,I_3)$ are independent across sample index $n$.
A classifier with known candidate parameters and equal initial probabilities
accumulates their exact Gaussian log likelihoods. Its probability after $n$
samples is
\begin{equation}
 \pi_k^{(n)}\propto\tfrac12\prod_{b=1}^{n}
       \Normal(o^{(b)};0,H\Sigma_k^0H^\top+\sigma^2I_3).
 \label{eq:fixed-classifier}
\end{equation}
The distinct symbol $\pi$ denotes the fixed classifier's probabilities; it
performs neither social updating nor parameter learning.

There are 200 independent sequences per generating candidate. Noise standard
deviations are $0$, $0.5$, $1$, and $2$; reported sample counts are
$1$, $5$, $20$, $100$, and $500$. Noise conditions share underlying state and
standard-noise draws, and sample counts use nested prefixes. Accuracy assigns
half credit to posterior ties and weights the two generating candidates equally.
Bootstrap intervals resample seed indices while keeping the two generating
candidates paired.

\begin{figure}[!htbp]
 \centering\includegraphics[width=0.79\linewidth]{fig_identifiability.pdf}
 \caption{Dependence-based discrimination. (a,b) Observable correlation matrices
 before measurement noise, with matched means and marginal variances.
 (c) Accuracy of the fixed-candidate classifier, with 200 sequences per generating
 candidate. $\sigma$ denotes measurement-noise standard deviation. Bands give
 pointwise 95\% seed-bootstrap intervals. The marginal-only classifier remains
 at chance. Sample counts are nested and noise conditions share underlying draws.}
 \label{fig:dependence}
\end{figure}

At $\sigma=1$, accuracy is 83.75\% after 20 samples and 99.75\% after 100;
at $\sigma=2$, it is 63.5\%, 82.5\%, and 97\% after 20, 100, and 500 samples.
Product-of-marginals and duplicate-candidate controls retain $\pi_k^{(n)}=1/2$;
label exchange preserves the aligned posterior. Dependence therefore distinguishes
this pair. Coefficient signs and paths both contribute: even $(C,M)$ has opposite
covariance signs. The check concerns known-model discrimination, not adaptive
social inference or unique causal-graph identification.
\FloatBarrier
\label{endappendix}
\end{document}
